\chapter{改造目标与总体思路}\label{chap:goal}

这次课程设计的目标，是把 Unix V6++ 中原本偏向连续进程镜像的内存管理方式，往真正按页组织的虚拟存储器模型推进。原来的系统虽然已经打开了硬件分页，但是内核侧很多逻辑仍然更像是在维护一整段连续的用户空间：进程增长、复制和释放都比较粗粒度，页表更多只是最后映射时才会碰到的结构。

我这次主要做了几件事。第一，把用户物理内存拆成 4KB 页框，用位图记录哪些页框已经被占用；第二，让每个进程拥有自己的页目录和私有用户页表，而不是所有进程都依赖同一份运行时页表；第三，在 \texttt{fork()}、\texttt{exec()}、\texttt{brk()}、栈增长和 \texttt{exit()} 这些路径里都改成逐页维护映射；最后通过 \texttt{pageDemo} 验证父子进程可以拥有相同的虚拟地址布局，但写入互不污染。

整体设计可以概括成下面这条线：

\begin{enumerate}
  \item 物理内存按页框分配和释放，不再只按连续空闲区来理解；
  \item 每个物理页框增加引用计数，为共享页和后续写时复制留下接口；
  \item 每个进程保存自己的页目录，页目录中同时挂接公共页表、内核页表和进程私有页表；
  \item 上下文切换时同步刷新 \texttt{CR3}，让 CPU 真正切换到当前进程的地址空间；
  \item 运行时通过 \texttt{pageDemo} 检查代码段、数据段、堆、栈的页布局，以及 \texttt{fork()} 后父子进程的隔离效果。
\end{enumerate}

所以这次改造并不是单独换了一个分配器，而是把进程地址空间的组织方式整体改成了“虚拟页到物理页框”的模型。后面的章节会按这个顺序说明实现细节。
